\documentclass[12pt]{article}
\usepackage{amsmath}
\usepackage{amssymb}
\usepackage{geometry}
\usepackage{ulem}

\geometry{margin=2.5cm}

\begin{document}

\section{Proving Q}

\noindent In order to compute the risk-neutral probability q, we need to start with the Replicated riskless Payoff H, which has value $V_1 =$ $H_{\text{up}}$ or $H_{\text{down}}$ depending on the states. Then, we want to maximise the expected utility of the replicated portofolio. This will allow us to derive the results and find the risk-neutral probability q, which makes each states indiferent.

\begin{center}
\begin{equation*}
V_1 = \alpha S_1 + \beta \cdot (1+r) =
\left\{
\begin{array}{ll}
\alpha \cdot S_0 u + \beta B_0 = X_{\text{up}} \\[4pt]
\alpha \cdot S_0 d + \beta B_0 = X_{\text{down}}
\end{array}
\right.
\end{equation*}
\end{center}
\vspace{1em}

\noindent\textbf{\uline{We want to maximise $E[u(v_1)]$:}}

\vspace{0.5em}

\[
\max_{\alpha,\, \beta} \; E\!\left[u\!\left(\alpha_i S_1 + \beta_i R\right)\right]
\quad \text{s.t.} \quad
\alpha_i S_0 + \beta_i B_0 = V_0
\]
\vspace{0.2em}

\noindent We maximise the Expected Utility using the Lagrangian Function:

\begin{align*}
\mathcal{L}(\alpha, \beta, \lambda)
  &= E\!\left[u\!\left(\alpha_i S_1 + \beta_i B_1\right)\right]
     - \lambda\!\left(\alpha_i S_0 + \beta_i B_0 - V_0\right) \\[6pt]
  &= p \cdot u(X_u) + (1-p) \cdot u(X_d)
     - \lambda\!\left(\alpha_i S_0 + \beta_i B_0 - V_0\right) \\[6pt]
  &= p \cdot u(\alpha S_{u} + \beta B_0)
     + (1-p) \cdot u(\alpha S_d + \beta B_0)
     - \lambda\!\left(\alpha_i S_0 + \beta B_0 - V_0\right)
\end{align*}

\vspace{1.5em}

\begin{align*}
\mathcal{L}_\alpha
  &= p \cdot u'(X_u) \cdot S_u + (1-p) \cdot u'(X_d) \cdot S_d
     - \lambda \cdot S_0 = 0
  \tag{I} \\[4pt]
  &= E\!\left[u'(V_1) \cdot S_1\right] = \lambda S_0
\end{align*}

\[
\longrightarrow \quad
\lambda = \frac{E\!\left[u'(V_1) \cdot S_1\right]}{S_0}
\]

\[
S_0 = \frac{1}{\lambda} \cdot E\!\left[u'(V_1) \cdot S_1\right]
\]

\vspace{2em}

\begin{align*}
\mathcal{L}_\beta
  &= p \cdot u'(X_u) \cdot B_1 + (1-p) \cdot u'(X_d) \cdot B_1
     - \lambda \cdot B_0 = 0
  \tag{II} \\[4pt]
  &= E\!\left[u'(V_1) \cdot B_1\right] = \lambda B_0
\end{align*}

\[
\longrightarrow \quad
\lambda = \frac{E\!\left[u'(V_1) \cdot B_1\right]}{B_0}
\]

\vspace{2em}

\noindent We assume that $B_0 = 1$ and $B_1 = (1+r)$, we thus have:

\[
\frac{1}{\lambda}
  = \frac{1}{(1+r) \cdot E\!\left[u'(V_1)\right]}
\]
\vspace{2em}


\noindent$\longrightarrow$ \underline{Substitute II in I:}
\[
S_0 = \frac{1}{1+r} \cdot \frac{E\!\left[u'(V_1) \cdot S_1\right]}{E\!\left[u'(V_1)\right]}
\]

\vspace{1.5em}

\noindent \textbf{\underline{We define the variable $Z$}}, which represents the marginal utility of $V_1$ divided by the
expected value of the marginal utility of $V_1$:
\[
Z = \frac{u'(V_1)}{E^P\!\left[u'(V_1)\right]}
\]
\vspace{0.1em}

\noindent Since $V_1$ can take two values, $Z$ takes also two values:
\[
Z_{\text{up}} =
  \frac{u'(X_{\text{up}})}
       {p \cdot u'(X_{\text{up}}) + (1-p) \cdot u'(X_{\text{down}})}
\qquad\qquad
Z_{\text{down}} =
  \frac{u'(X_{\text{down}})}
       {p \cdot u'(X_{\text{up}}) + (1-p) \cdot u'(X_{\text{down}})}
\]

\vspace{1.5em}

\noindent\uline{$Z$ has two properties:}

\begin{enumerate}
  \item Since we assume increasing utility $\bigl(u'(V_1) > 0\bigr)$:
        \quad $Z > 0$ \quad $\forall\, \mathbb{R}$

  \vspace{0.6em}

  \item The expected value of $Z$ is equal to 1:
        \[
          E^P[Z]
          = E^P\!\left[\frac{u'(V_1)}{E^P\!\left[u'(V_1)\right]}\right]
          = \frac{E^P\!\left[u'(V_1)\right]}{E^P\!\left[u'(V_1)\right]}
          = 1
        \]
\end{enumerate}

\vspace{1.5em}

\noindent Therefore, we can rewrite equation iii as:
\begin{align*}
              & S_0 = \frac{1}{1+r} \cdot E^P\!\left[Z \cdot S_1\right] \\[0.5pt]
\intertext{\footnotesize Where: \quad $E^P\!\left[Z \cdot S_1\right] = p \cdot Z_{\text{up}} \cdot S_{0,u} + (1-p) \cdot Z_{\text{down}} \cdot S_{0,d}$}
\\[-8pt]
\longrightarrow \quad & S_0 = \frac{1}{1+r} \cdot p \cdot Z_{\text{up}} \cdot S_{0,u} + (1-p) \cdot Z_{\text{down}} \cdot S_{0,d}
\end{align*}

\vspace{1.5em}

\noindent By dividing both sides by $S_0$, we get:
\[
1 + r = p \cdot Z_{\text{up}} \cdot u + (1-p) \cdot Z_{\text{down}} \cdot d
\]

\vspace{1.5em}

\noindent We define $q$ and $(1-q)$ as:
\[
q = p \cdot Z_{\text{up}}
  = p \cdot \frac{u'(X_{\text{up}})}{E\!\left[u'(V_1)\right]}
\qquad \text{and} \qquad
1-q = (1-p) \cdot Z_{\text{down}}
       = (1-p) \cdot \frac{u'(X_{\text{down}})}{E\!\left[u'(V_1)\right]}
\]
\vspace{1.5em}

We thus have: 
\begin{align*}
\longrightarrow \quad & 1 + r = q \cdot u + (1-q) \cdot d \\[6pt]
                      & 1 + r = q(u - d) + d
\end{align*}

\[
\boxed{ q = \frac{(1+r) - d}{u - d} }
\]

\noindent -> still adding text
\end{document}